\subsection{An isolated Bernoulli test for an even class-group component}
\label{sec:isolated-reflection}

\begin{lemma}[K17-REFLECTOR]\label{lem:isolated-reflection}
Let $p$ be an odd prime and let $e$ be even with $2\le e\le p-3$.
Write $A_p$ for the $p$-primary ideal class group of
$F=\mathbf Q(\mu_p)$, with $A_p^{[i]}$ denoting its
$\omega^i$-eigenspace. Then
\[
 p\nmid B_{p-1-e}
 \quad\Longrightarrow\quad A_p^{[1+e]}=0
 \quad\Longrightarrow\quad A_p^{[-e]}=0.
 \tag{RB.1}
\]
Here divisibility of a rational Bernoulli number means divisibility of
its numerator in lowest terms. In particular, the left-hand condition
implies $K_{2e}(\mathbf Z)\{p\}=0$ by (K2).
\end{lemma}

\begin{proof}
Put $n=p-1-e$. This is even and satisfies $2\le n\le p-3$.
Herbrand's implication, in the precise form of
\cite[Main Theorem (1.1), p.~151]{Ribet1976}, says that
$(A_p/pA_p)^{[1-n]}\ne0$ implies $p\mid B_n$.
Since $1-n\equiv1+e\pmod{p-1}$ and a finite $p$-group is zero
if its quotient by $p$ is zero, this proves the first implication.
Only the Herbrand direction is used; Ribet's converse is unnecessary.

For completeness, we prove the needed direction of reflection.
It also appears in the full proof of
\cite[Theorem~4.4.2, pp.~111--113, and Corollary~4.4.4,
p.~113]{K17SharifiReflection}.
Set $i=-e$ and $j=1-i=1+e$ modulo $p-1$.
Then $i$ is even and nonzero, while $j$ is odd and distinct from $1$.
Class field theory identifies
\[
 B=A_p^{[i]}/pA_p^{[i]}
\]
with the Galois group of a finite elementary abelian extension $L/F$
unramified at every finite prime. The corresponding quotient is stable
under $\Delta=\operatorname{Gal}(F/\mathbf Q)$, so $L$ is Galois
over $\mathbf Q$ and this identification is $\Delta$-equivariant.

Let
\[
 V=\{a\in F^\times:a^{1/p}\in L\}/F^{\times p}
\]
be its Kummer radical. The perfect, equivariant Kummer pairing
$B\times V\to\mu_p$ gives
$V=V^{[j]}$ and $\dim_{\mathbf F_p}V=\dim_{\mathbf F_p}B$.
For any representative $a$ of $V$ and any finite valuation $v$ of $F$,
choose a valuation above $v$ in $L$. Since its ramification index is one,
the equation $(a^{1/p})^p=a$ shows that $p$ divides $v(a)$.
Consequently $(a)=\mathfrak a^p$ for a unique fractional ideal
$\mathfrak a$. The assignment
\[
             V\longrightarrow A_p[p],\qquad
             a\longmapsto[\mathfrak a]
\]
is well-defined and $\Delta$-equivariant. Its kernel is contained in
the image of $\mathcal O_F^\times/\mathcal O_F^{\times p}$:
if $\mathfrak a=(b)$, then $a/b^p$ is a unit.

The odd part of this unit quotient has only character $\omega$.
Indeed, for a unit $u$, every complex conjugate of $u/\overline u$
has absolute value one, so this algebraic integer unit is a root of
unity. Moreover
\[
 u^2=(u/\overline u)(u\overline u).
\]
Thus the quotient of the unit group by roots of unity and real units
is killed by two. Modulo $p$th powers, the real units contribute only
even characters, and the $p$-primary roots of unity in $F$ contribute
exactly the character $\omega$. Since $j\ne1$, the above map is
injective on $V^{[j]}$. Hence
\[
 \dim_{\mathbf F_p}(A_p/pA_p)^{[i]}
 \le \dim_{\mathbf F_p}A_p^{[j]}[p].
\]
The right-hand side is zero if $A_p^{[j]}=0$, and this proves the
second implication. Finally $e=2k$ for a positive integer $k$, so (K2)
identifies the vanishing of $A_p^{[-e]}$ with that of
$K_{2e}(\mathbf Z)\{p\}$.
\end{proof}

\paragraph{Exact specialization of the remaining degree-twenty prime.}
For
\[
                 p=9496386083,\qquad e=10,
\]
the single Bernoulli index needed by (RB.1) is
$n=9496386072$. The reflected component is $[11]$, since
$1-n\equiv11\pmod{p-1}$. No verification of all even components at
this prime, and no assumption of Vandiver there, is required. The
separate exact Bernoulli certificate supplies the numerical residue;
the present lemma specifies its full arithmetic implication and the
character conventions.
